\chapter{Determinants}

\section{Definition by cofactor expansion}

For $2\times2$ matrices we already met $\det\begin{bmatrix}a&b\\c&d\end{bmatrix}=ad-bc$, which decides
invertibility. We extend this recursively.

\begin{definition}[Minor, cofactor, determinant]
Let $A$ be $n\times n$ with $n\ge2$. Let $A(i,j)$ be the $(n-1)\times(n-1)$ matrix obtained by deleting
row $i$ and column $j$ of $A$. The \emph{cofactor} of $a_{ij}$ is
\[
  C_{ij}=(-1)^{i+j}\det A(i,j).
\]
The \emph{determinant} is $\det A = a_{11}C_{11}+a_{12}C_{12}+\dots+a_{1n}C_{1n}$ (expansion along the first
row), with $\det[a]=a$ for $1\times1$ matrices. We also write $\abs{A}$ for $\det A$.
\end{definition}

The signs $(-1)^{i+j}$ follow a checkerboard pattern starting with $+$ in the top-left corner.

\begin{theorem}[Cofactor expansion along any row or column]
For any $i$ and any $j$,
\[
  \det A=\sum_{k=1}^n a_{ik}C_{ik}=\sum_{k=1}^n a_{kj}C_{kj}.
\]
\end{theorem}

The proof is a (tedious) induction on $n$. Practically: expand along a row or column with many zeros.

\begin{example}
For $A=\begin{bmatrix}1&2&3\\2&5&3\\1&0&8\end{bmatrix}$, along the first row:
\[
\det A = 1\begin{vmatrix}5&3\\0&8\end{vmatrix}-2\begin{vmatrix}2&3\\1&8\end{vmatrix}+3\begin{vmatrix}2&5\\1&0\end{vmatrix}
=1(40)-2(16-3)+3(0-5)=40-26-15=-1.
\]
Along the second column (which has a zero):
$\det A = -2\begin{vmatrix}2&3\\1&8\end{vmatrix}+5\begin{vmatrix}1&3\\1&8\end{vmatrix}-0
= -2(13)+5(5) = -1$. Same answer.
\end{example}

\begin{theorem}
If $A$ is upper or lower triangular, $\det A$ is the product of the diagonal entries. In particular
$\det I=1$. If $A$ has a row or column of zeros, $\det A=0$.
\end{theorem}
\begin{proof}
For lower triangular, expand along the first row: only $a_{11}C_{11}$ survives, and $A(1,1)$ is again lower
triangular; induct. (Upper triangular: expand along the first column.)
\end{proof}

\section{Determinants and row operations}

\begin{theorem}[Effect of row operations]
Let $A$ be $n\times n$.
\begin{enumerate}
  \item If $B$ is obtained by multiplying one row of $A$ by $c$, then $\det B=c\det A$.
  \item If $B$ is obtained by swapping two rows, then $\det B=-\det A$.
  \item If $B$ is obtained by adding a multiple of one row to another, then $\det B=\det A$.
\end{enumerate}
\end{theorem}
\begin{proof}[Proof sketch]
(1) Expand along the scaled row. (2) Induction: for $n\ge3$, expand along a row not involved in the swap; each
minor has two rows swapped. The $2\times2$ case is direct. Consequently a matrix with two equal rows has
determinant $0$ (swapping them gives $\det A=-\det A$). (3) Expanding $B$ along row $i$ (where
$R_i\to R_i+cR_j$) gives $\det A + c\det A'$, where $A'$ has row $j$ appearing twice, so $\det A'=0$.
\end{proof}

\begin{important}[Strategy]
Row reduce to upper triangular form, tracking swaps (each flips the sign) and scalings (each multiplies
by the scalar), then multiply the diagonal. Because $\det A^T=\det A$, column operations behave the same way.
\end{important}

\begin{example}
Compute $\det B$ for $B=\begin{bmatrix}1&2&0&1\\2&5&1&3\\0&1&3&2\\1&2&1&4\end{bmatrix}$.
\[
\begin{vmatrix}1&2&0&1\\2&5&1&3\\0&1&3&2\\1&2&1&4\end{vmatrix}
\overset{\substack{R_2-2R_1\\R_4-R_1}}{=}
\begin{vmatrix}1&2&0&1\\0&1&1&1\\0&1&3&2\\0&0&1&3\end{vmatrix}
\overset{R_3-R_2}{=}
\begin{vmatrix}1&2&0&1\\0&1&1&1\\0&0&2&1\\0&0&1&3\end{vmatrix}
\]
\[
\overset{R_3\leftrightarrow R_4}{=}
-\begin{vmatrix}1&2&0&1\\0&1&1&1\\0&0&1&3\\0&0&2&1\end{vmatrix}
\overset{R_4-2R_3}{=}
-\begin{vmatrix}1&2&0&1\\0&1&1&1\\0&0&1&3\\0&0&0&-5\end{vmatrix}.
\]
So $\det B=-(1\cdot1\cdot1\cdot(-5))=5$.
\end{example}

\begin{theorem}
$A$ is invertible if and only if $\det A\ne0$.
\end{theorem}
\begin{proof}
Row reduce $A$ to its RREF $R$; each operation multiplies the determinant by a non-zero number, so
$\det A\neq0$ iff $\det R\ne0$. For square RREF, $R=I$ (determinant $1$) if $\rank A=n$, and otherwise $R$ has
a zero row (determinant $0$).
\end{proof}

\section{Products and transposes}

\begin{theorem}
For $n\times n$ matrices $A,B$:
\begin{enumerate}
  \item $\det(EA)=\det E\det A$ for any elementary matrix $E$;
  \item $\det(AB)=\det A\det B$;
  \item $\det(A^T)=\det A$;
  \item if $A$ is invertible, $\det(A^{-1})=\dfrac{1}{\det A}$;
  \item $\det(cA)=c^n\det A$.
\end{enumerate}
\end{theorem}
\begin{proof}[Proof sketch]
(1) is the row-operation theorem, since $\det E$ is $c$, $-1$ or $1$ respectively. (2) If $A$ is invertible,
write $A=E_1\cdots E_k$ and apply (1) repeatedly. If not, $AB$ is not invertible either (if $AB$ had inverse
$C$ then $A(BC)=I$ would make $A$ invertible), so both sides are $0$. (4) follows from
$\det A\det A^{-1}=\det I=1$. (5): $cA$ scales each of $n$ rows by $c$.
\end{proof}

\begin{important}[Pitfall]
$\det(A+B)\neq\det A+\det B$ in general, and $\det(cA)=c^n\det A$, \emph{not} $c\det A$.
For example $\det(2I_3)=8$.
\end{important}

\section{The adjugate and a formula for the inverse}

\begin{definition}
The \emph{cofactor matrix} of $A$ is $\cof A=[C_{ij}]$, and the \emph{adjugate} is
$\adj A = (\cof A)^T$, i.e.\ $(\adj A)_{ij}=C_{ji}$.
\end{definition}

\begin{theorem}
$A\,\adj A=(\det A)I=\adj A\,A$. Hence if $\det A\ne0$, $A^{-1}=\dfrac{1}{\det A}\adj A$.
\end{theorem}
\begin{proof}
$(A\adj A)_{ij}=\sum_k a_{ik}C_{jk}$. If $i=j$ this is the expansion of $\det A$ along row $i$. If $i\ne j$ it is
the expansion along row $j$ of the matrix obtained from $A$ by replacing row $j$ with row $i$; that matrix has
two equal rows, so the sum is $0$ (``expansion by alien cofactors'').
\end{proof}

\begin{example}
For $A=\begin{bmatrix}1&2&3\\2&5&3\\1&0&8\end{bmatrix}$ (with $\det A=-1$), the cofactors are
\[
\begin{aligned}
C_{11}&=40, & C_{12}&=-(16-3)=-13, & C_{13}&=0-5=-5,\\
C_{21}&=-(16-0)=-16, & C_{22}&=8-3=5, & C_{23}&=-(0-2)=2,\\
C_{31}&=6-15=-9, & C_{32}&=-(3-6)=3, & C_{33}&=5-4=1.
\end{aligned}
\]
So $\adj A=\begin{bmatrix}40&-16&-9\\-13&5&3\\-5&2&1\end{bmatrix}$ and
$A^{-1}=\frac{1}{-1}\adj A=\begin{bmatrix}-40&16&9\\13&-5&-3\\5&-2&-1\end{bmatrix}$, agreeing with the
row-reduction computation of Chapter 7.
\end{example}

For a $2\times2$ matrix the adjugate is $\begin{bmatrix}d&-b\\-c&a\end{bmatrix}$, recovering the familiar formula.

\section{Cramer's Rule}

\begin{theorem}[Cramer's Rule]
If $A$ is invertible, the unique solution of $A\vec x=\vec b$ is given by
\[
  x_i=\frac{\det A_i}{\det A},\qquad i=1,\dots,n,
\]
where $A_i$ is $A$ with its $i$-th column replaced by $\vec b$.
\end{theorem}
\begin{proof}
$\vec x=A^{-1}\vec b=\frac{1}{\det A}\adj A\,\vec b$, so
$x_i=\frac{1}{\det A}\sum_k C_{ki}b_k$, and $\sum_kb_kC_{ki}$ is the expansion of $\det A_i$ along column $i$.
\end{proof}

\begin{example}
Solve $3x_1+2x_2=7$, $x_1+4x_2=9$. Here $\det A=12-2=10$,
\[
x_1=\frac{1}{10}\begin{vmatrix}7&2\\9&4\end{vmatrix}=\frac{28-18}{10}=1,\qquad
x_2=\frac{1}{10}\begin{vmatrix}3&7\\1&9\end{vmatrix}=\frac{27-7}{10}=2.
\]
Check: $3+4=7$, $1+8=9$.
\end{example}

Cramer's rule is mostly of theoretical value (e.g.\ it shows the solution depends continuously on the
entries); for computation, row reduction is far cheaper.

\section{Area and volume}

\begin{theorem}
If $\vec u,\vec v\in\R^2$, the parallelogram they span has area $\abs{\det[\vec u\ \vec v]}$. If
$\vec u,\vec v,\vec w\in\R^3$, the parallelepiped they span has volume $\abs{\det[\vec u\ \vec v\ \vec w]}
=\abs{\vec u\cdot(\vec v\times\vec w)}$. More generally, a linear map $L:\R^n\to\R^n$ multiplies
$n$-dimensional volumes by $\abs{\det[L]}$.
\end{theorem}
\begin{proof}[Idea]
Area/volume is unchanged by shears (adding a multiple of one column to another), scales by $\abs{c}$ when one
edge is scaled by $c$, and is $1$ for the unit square/cube; the determinant behaves the same way (up to sign).
\end{proof}

\begin{example}
The parallelogram spanned by $(3,1)$ and $(1,2)$ has area $\abs{\det\begin{bmatrix}3&1\\1&2\end{bmatrix}}=\abs{6-1}=5$.
The parallelepiped spanned by $(1,2,0)$, $(0,1,3)$, $(1,0,1)$ has volume
\[
\abs{\det\begin{bmatrix}1&0&1\\2&1&0\\0&3&1\end{bmatrix}}
=\abs{1(1\cdot1-0\cdot3)-0+1(2\cdot3-1\cdot0)}=\abs{1+6}=7.
\]
The sign of the determinant records orientation: a negative determinant means the map reverses
orientation (like a reflection).
\end{example}

\section*{Summary}
\begin{itemize}
  \item $\det A$ by cofactor expansion along any row or column; triangular: product of diagonal.
  \item Row operations: scaling by $c$ multiplies $\det$ by $c$; swap negates; adding a multiple of a row
        changes nothing.
  \item $\det(AB)=\det A\det B$, $\det A^T=\det A$, $\det A^{-1}=1/\det A$, $\det(cA)=c^n\det A$.
  \item $A$ invertible $\iff\det A\ne0$; $A^{-1}=\frac{1}{\det A}\adj A$; Cramer's rule
        $x_i=\det A_i/\det A$.
  \item $\abs{\det}$ = area/volume scaling factor.
\end{itemize}
